%==========================================================================
\section{Discussion}\label{sec:discussion}

The measurements separate three regimes.

\textbf{Where the FPGA wins.} Circuits whose passes carry few gates per amplitude, such as GHZ, are bound by memory bandwidth. Here 16 engines streaming their own channels outrun a 128-core server from \NumXplatGhzCrossN{} qubits upwards, even though the FPGA time includes returning the state to the host and sampling, work the CPU runs do not do. Against the baseline kernel on the same card the redesign is faster and more energy-efficient at every size, because the baseline used a small fraction of the available bandwidth (Section~\ref{sec:baseline}). The kernel time is also predictable from run to run, although the current host path makes the end-to-end time less predictable than on the CPU (Section~\ref{subsec:res-jitter}).

\textbf{Where it ties or loses.} Circuits with many gates per pass (QFT, random circuits, the VQE ansatz) are bound by the arithmetic of the engines, whose peak of \NumFlagComputeRoof~GFLOP/s is set by the DSP budget of 16 engines with four pair updates per cycle (Fig.~\ref{fig:roofline}). For these circuits the 128-core server remains faster at 30 qubits, by a factor of \NumXplatCpuFasterOtherThirtyMin to \NumXplatCpuFasterOtherThirtyMax; its fastest configurations for them fuse consecutive gates into wider matrices, which raises the arithmetic intensity on the CPU side as well. Small circuits are dominated by a fixed host and launch overhead of a few milliseconds, against tens of microseconds on the CPU, so the card is not suited to workloads of many small independent circuits unless the host path is moved to compiled code or several circuits are batched in one launch. The A100 GPU is faster than both the card and the best CPU configuration for QFT and RQC at 30 qubits (\NumXplatBestGpuRatioQftThirty$\times$ the card for QFT); for VQE it is about as fast as the best CPU configuration, and for GHZ its fastest configuration and the card are about equal. Over the four families the best CPU configuration needs \NumXplatGpuOverCpuThirtyMin~to \NumXplatGpuOverCpuThirtyMax~times the GPU time at 30 qubits.

\textbf{What would change the picture.} The remaining overheads at 30 qubits are exchange (Fig.~\ref{fig:breakdown}) and the arithmetic bound. Both are design choices rather than limits of the card: the measured 32-channel ceiling is twice what 16 engines can use, and the layout ablation (Table~\ref{tab:ablation}) shows that the placement of qubits changes the number of remaps, so a better placement policy would reduce exchange. Precision is a second axis: the engine computes in complex64, while Qiskit Aer in complex128 with fusion is also faster than the FPGA for QFT, RQC and VQE at 30 qubits (Table~\ref{tab:xplat}).


\subsection{Future work: several cards}\label{subsec:future-multifpga}

A single U55C limits the statevector to 30 qubits in complex64. The architecture extends to several cards in the same way it uses several pseudo-channels on one card: the card index becomes the top qubits of the physical index, gates on those qubits require an exchange between cards, and the host scheduler can minimise such exchanges with the same remapping policy that it uses for bank qubits. Because PCIe offers far less bandwidth than HBM (Section~\ref{subsec:res-ceilings}), the exchange between cards, not the engines, is expected to dominate. A distributed statevector architecture on four Alveo U55C cards is the subject of a separate paper~\cite{ali2026multifpga}; this paper does not report multi-card measurements.
